\begin{tikzpicture}
	
	%%% prediction accuracy of synthetic enhancer strength with different inputs %%%
	%% explain input data
	\coordinate (scheme) at (0, 0);
	
	\foreach \x in {1, 2, 3}{
		\node[anchor = west, node font = \figsmall, yshift = -\x * .75cm] (seq\x) at (scheme) {seq\x};
	}
	
	\distance{seq3.east}{\threecolumnwidth, 0}
	\pgfmathsetlength{\templength}{(\xdistance - 2\columnsep) / 4}
	
	\pgfmathsetseed{928}
	
	\colorlet{seq1}{gray!75!white}
	\colorlet{seq2}{gray}
	\colorlet{seq3}{gray!75!black}
	
	% sequence- and position-specific effect
	\node[anchor = north west, node font = \figsmall\bfseries, text width = \templength, align = center, inner xsep = 0pt, shift = {(.5\columnsep, -.5\columnsep)}] (seq+pos) at (scheme -| seq3.east) {seq+pos};
	
	\foreach \x in {1, 2, 3}{
		\draw[ultra thick, seq\x] (seq\x -| seq+pos.west) -- ++(\templength, 0) coordinate[pos = .165] (p1) coordinate[pos = .5] (p2) coordinate[pos = .835] (p3);
		\foreach \y in {1, 2, 3}{
			\draw[->] (p\y) +(0, .05) -- ++(0, -.35) node[anchor = north, node font = \figtiny] {%
				\pgfmathparse{rnd}%
				\expandafter\xdef\csname tmp\x\y\endcsname{\pgfmathresult}%
				\pgfmathprintnumber[fixed, zerofill, precision = 1]{\csname tmp\x\y\endcsname}%
			};
		}
	}
	
	% sequence-specific effect
	\node[anchor = base west, node font = \figsmall\bfseries, text width = \templength, align = center, inner xsep = 0pt, xshift = .5\columnsep] (seq only) at (seq+pos.base east) {seq only};
	
	\foreach \x in {1, 2, 3}{
		\draw[ultra thick, seq\x] (seq\x -| seq only.west) -- ++(\templength, 0) coordinate[pos = .165] (p1) coordinate[pos = .5] (p2) coordinate[pos = .835] (p3);
		\coordinate[yshift = -.35cm] (e\x) at (p2);
		\foreach \y in {1, 2, 3}{
			\draw[->] (p\y) +(0, .05) -- ++(0, -.05) to[out = -90, in = 90] ($(e\x) + (0, .05)$) -- (e\x);
		}
		\node[anchor = north, node font = \figtiny] at (e\x) {%
			\pgfmathparse{(\csname tmp\x1\endcsname + \csname tmp\x2\endcsname + \csname tmp\x3\endcsname) / 3}%
			\pgfmathprintnumber[fixed, zerofill, precision = 1]{\pgfmathresult}%
		};
	}
	
	% position-specific effect
	\node[anchor = base west, node font = \figsmall\bfseries, text width = \templength, align = center, inner xsep = 0pt, xshift = .5\columnsep] (pos only) at (seq only.base east) {pos only};
	
	\foreach \x in {1, 2, 3}{
		\draw[ultra thick, seq\x] (seq\x -| pos only.west) -- ++(\templength, 0) coordinate[pos = .165] (p1) coordinate[pos = .5] (p2) coordinate[pos = .835] (p3);
	}
	\foreach \y in {1, 2, 3}{
		\draw[->] (seq1 -| p\y) +(0, .05) -- (e3 -| p\y) node[anchor = north, node font = \figtiny] {%
			\pgfmathparse{(\csname tmp1\y\endcsname + \csname tmp2\y\endcsname + \csname tmp3\y\endcsname) / 3}%
			\pgfmathprintnumber[fixed, zerofill, precision = 1]{\pgfmathresult}%
		};
	}
	
	% mean effect
	\node[anchor = base west, node font = \figsmall\bfseries, text width = \templength, align = center, inner xsep = 0pt, xshift = .5\columnsep] (mean) at (pos only.base east) {mean};
	
	\foreach \x in {1, 2, 3}{
		\draw[ultra thick, seq\x] (seq\x -| mean.west) -- ++(\templength, 0) coordinate[pos = .165] (p1) coordinate[pos = .5] (p2) coordinate[pos = .835] (p3);
	}
	\foreach \y in {1, 2, 3}{
		\draw[->] (seq1 -| p\y) +(0, .05) -- (seq3 -| p\y) -- ++(0, -.05) to[out = -90, in = 90] ($(p2 |- e3) + (0, .05)$) -- (p2 |- e3);
	}
	\node[anchor = north, node font = \figtiny] (effect) at (p2 |- e3) {%
		\pgfmathparse{%
			(%
				\csname tmp11\endcsname + \csname tmp12\endcsname + \csname tmp13\endcsname +%
				\csname tmp21\endcsname + \csname tmp22\endcsname + \csname tmp23\endcsname +%
				\csname tmp31\endcsname + \csname tmp32\endcsname + \csname tmp33\endcsname
			) / 9%
		}%
		\pgfmathprintnumber[fixed, zerofill, precision = 1]{\pgfmathresult}%
	};
	
	% ACR lib
	\node[anchor = north east, node font = \figsmall\bfseries, text width = \templength, align = right, yshift = -\columnsep] (ACR lib) at (seq+pos.east |- effect.south) {ACR lib:};
	\node[anchor = base west, node font = \figsmall, text width = 3\templength, align = left] at (ACR lib.base east) {transcription factor effects from ACR library};
	
	% random
	\node[anchor = north east, node font = \figsmall\bfseries, text width = \templength, align = right, yshift = -3\columnsep] (random) at (seq+pos.east |- effect.south) {random:};
	\node[anchor = base west, node font = \figsmall] at (random.base east) {random permutation of \textbf{seq+pos}};
	
	
	%% plot prediction accuracy
	\coordinate (accuracy) at (scheme -| \textwidth - \twothirdcolumnwidth, 0);
	
	\expandafter\def\csname id+pos\endcsname{\vphantom{ly}seq+pos}
	\expandafter\def\csname no_pos\endcsname{\vphantom{ly}seq only}
	\expandafter\def\csname no_id\endcsname{\vphantom{ly}pos only}
	\expandafter\def\csname mean\endcsname{\vphantom{ly}mean}
	\expandafter\def\csname tamsACR\endcsname{\vphantom{ly}ACR lib}
	\expandafter\def\csname random\endcsname{\vphantom{ly}random}
	
	\begin{pggfplot}[%
		at = {(accuracy)},
		total width = \twothirdcolumnwidth,
		xlabel = method to obtain transcription factor effects,
		ylabel = prediction accuracy ($R^2$),
		ymin = 0,
		enlarge x limits = 0,
		enlarge y limits = {lower = 0},
		title is csname,
		title color from name,
		xticklabels are csnames,
		xticklabel style = {rotate = 45, anchor = north east, inner sep = 0.236em, text depth = -2pt},
		tight y
	]{synEnh_pred_accuracy}
	
		\pggf_bar(
			ybar,
			color from column name,
			draw = black
		)
	
	\end{pggfplot}
	
	
	%%% predict synthetic enhancer strength %%%
	\coordinate[yshift = -\columnsep] (correlation) at (scheme |- xlabel.south);
	
	\begin{pggfplot}[%
		at = {(correlation)},
		total width = \textwidth,
		xlabel = {\textbf{measurement}: \enrichment},
		ylabel = {\textbf{prediction (seq+pos)}:\\\enrichment},
		ylabel add to width = 3pt,
		facet 1/.append style = {
			pggf colorbar = {%
				title = count,
				log2,
				horizontal,
				ytick = {0, 2, ..., 15}
			}
		},
		title is csname,
		title color from name,
		xytick = {-12, -10, ..., 12},
	]{synEnh_prediction}
	
		\pggf_annotate(diagonal)
	
		\pggf_hexbin(bins = 50)
	
		\pggf_stats(stats = {n,spearman,rsquare})
	
	\end{pggfplot}
	
	
	%%% subfigure labels %%%
	\subfiglabel{scheme}
	\subfiglabel{accuracy}
	\subfiglabel{correlation}

\end{tikzpicture}